\chapter{Exercise Solutions}
\label{app:solutions}

This appendix provides worked solutions to the exercises at the end of
each chapter. The solutions give the approach, the key code or formula,
and the expected result. Where an exercise asks for experimentation, we
report the qualitative finding rather than a specific number---the
number depends on the random seed and data window.

\begin{keyinsight}
The exercises are designed to be reproducible from the bundled data
and the companion crates. Every solution below can be verified by
running the corresponding example in \crate{quant-lib} with minor
modifications. The goal is understanding, not memorisation.
\end{keyinsight}

\section*{Runnable Code}
\label{app:solutions:code}

Every exercise in this appendix has a complete, runnable Rust
solution in the \crate{quant-lib} example suite. The 70 solutions
span chapters 1--14 and live in
\texttt{crates/quant-lab/crates/quant-lib/examples/solutions/}, one
file per chapter:

\begin{center}
\begin{tabular}{lll}
\textbf{Chapter} & \textbf{Topic} & \textbf{Example target} \\
1  & Credit Card Fraud       & \texttt{solutions-ch01\_fraud\_solutions} \\
2  & Loan Default            & \texttt{solutions-ch02\_loan\_solutions} \\
3  & Market Data             & \texttt{solutions-ch03\_stocks\_solutions} \\
4  & Returns \& Volatility   & \texttt{solutions-ch04\_returns\_solutions} \\
5  & Backtesting             & \texttt{solutions-ch05\_backtest\_solutions} \\
6  & Statistical Foundations & \texttt{solutions-ch06\_moments\_solutions} \\
7  & Time Series             & \texttt{solutions-ch07\_timeseries\_solutions} \\
8  & Volatility Models       & \texttt{solutions-ch08\_volatility\_solutions} \\
9  & Stochastic Processes    & \texttt{solutions-ch09\_stochastic\_solutions} \\
10 & Options Pricing         & \texttt{solutions-ch10\_options\_solutions} \\
11 & Portfolio Optimization  & \texttt{solutions-ch11\_portfolio\_solutions} \\
12 & Factor Models           & \texttt{solutions-ch12\_factors\_solutions} \\
13 & Market Microstructure   & \texttt{solutions-ch13\_microstructure\_solutions} \\
14 & AFML Backtesting        & \texttt{solutions-ch14\_afml\_solutions} \\
\end{tabular}
\end{center}

Run any solution and inspect its output:
\begin{lstlisting}[language=bash]
cargo run -p quant-lib --example solutions-ch01_fraud_solutions
\end{lstlisting}
Each file also carries one \texttt{\#\[test\]} per exercise (70 tests
total), runnable with:
\begin{lstlisting}[language=bash]
cargo test -p quant-lib --examples
\end{lstlisting}

Five end-to-end real-world projects tie the chapters together and live
in \texttt{crates/quant-lab/crates/quant-lib/examples/projects/}:

\begin{center}
\begin{tabular}{lll}
\textbf{Project} & \textbf{Level} & \textbf{Example target} \\
1 & Cross-sectional momentum    & \texttt{projects-01\_momentum} \\
2 & Pairs trading (cointegration) & \texttt{projects-02\_pairs\_trading} \\
3 & Risk-parity portfolio        & \texttt{projects-03\_risk\_parity} \\
4 & Options market making        & \texttt{projects-04\_options\_mm} \\
5 & Full AFML pipeline           & \texttt{projects-05\_afml\_pipeline} \\
\end{tabular}
\end{center}

The project catalogue and per-project metrics are documented in
\texttt{examples/projects/README.md}. The prose in the remainder of
this appendix explains the approach, the key formula, and the
qualitative finding for each exercise; the runnable code is the
single source of truth for the implementation.

\section*{Chapter 1: Credit Card Fraud}

\paragraph{1. Threshold Sensitivity.}
Sweep the threshold over $\{2.0, 2.5, 3.0, 3.5, 4.0\}$, computing
precision and recall at each. Lowering the threshold raises recall
(more transactions flagged) but lowers precision (more false alarms).
The F1-maximising threshold is the one where the precision--recall
curve has its knee---typically around $3.0$--$3.5$ for the bundled
sample. The lesson: there is no ``correct'' threshold; the choice
depends on the cost of a false negative versus a false positive.

\paragraph{2. Feature Analysis.}
Modify \func{predict} to return the per-feature z-scores alongside
the boolean. The most discriminative feature is the one whose
z-score most consistently exceeds the threshold on fraud samples.
In the anonymised PCA dataset (\texttt{V1}--\texttt{V28}), features
\texttt{V14} and \texttt{V4} typically dominate---they carry the
highest mean $|z|$ on the fraud rows.

\paragraph{3. Train/Test Split.}
Split chronologically because fraud patterns evolve: a random split
leaks future information into training. The test-set performance is
almost always worse than training-set: precision drops by 10--20\%
because the z-score statistics are calibrated on the earlier period.
This is the simplest demonstration of why time-series data must
never be shuffled before splitting.

\paragraph{4. Amount Feature.}
Add \texttt{amount} as a z-scored feature. Performance barely moves
because the PCA features already absorb most of the signal; the
raw amount is weakly correlated with fraud after conditioning on the
anonymised features. The lesson: not every available feature adds
value---check before including.

\paragraph{5. Trait Extraction.}
The \texttt{AnomalyDetector} trait should require:
\begin{lstlisting}[language=Rust]
pub trait AnomalyDetector {
    fn fit(&mut self, x: &[Vec<f64>]);
    fn predict(&self, x: &[f64]) -> AnomalyResult;
}

pub struct AnomalyResult {
    pub is_anomaly: bool,
    pub score: f64,
    pub triggering_features: Vec<usize>,
}
\end{lstlisting}
\func{fit} learns the per-feature mean and std; \func{predict}
returns the z-score, the boolean, and which features exceeded the
threshold. This is the same shape as \func{BinaryClassifier} plus
\func{Scorer} from \crate{qf-02-loan}, generalised to multivariate
input.

\section*{Chapter 2: Loan Default}

\paragraph{1. Weight Sensitivity.}
Sweep each weight $\pm 50\%$ while holding the others fixed. The AUC
moves by at most a few points, indicating the linear scorer is
robust to moderate weight mis-specification. The most sensitive
feature is the one whose weight change moves AUC the most---typically
\texttt{debt\_to\_income} or \texttt{credit\_score}, depending on the
dataset.

\paragraph{2. Threshold Selection.}
The sigmoid output $\sigma(\mathbf{w}^\top \mathbf{x})$ is a
probability. Sweep the threshold from $0.01$ to $0.99$; the F1
maximum is typically near the base rate (the fraction of defaults
in the data), because that is where precision and recall balance
for a balanced classifier.

\paragraph{3. Feature Importance.}
Train a single-feature scorer for each feature and compute its AUC.
Rank by AUC. The ranking usually matches the weight magnitude from
the multivariate scorer, but not always: correlated features can
split importance, so a feature that is weak alone can be strong in
combination.

\paragraph{4. Cross-Validation.}
Split into 5 chronological folds. The AUC standard deviation across
folds is typically 0.02--0.05. A large std (relative to the mean)
indicates the model is sensitive to the training window---a sign of
instability, not robustness.

\paragraph{5. Calibration Plot.}
Bin predictions into deciles and plot the mean predicted probability
against the empirical default rate in each bin. A well-calibrated
model lies on the diagonal; a miscalibrated one is systematically
above (overconfident) or below (underconfident). The linear scorer
is usually underconfident at high probabilities because the sigmoid
saturates.

\section*{Chapter 3: Market Data}

\paragraph{1. Candlestick Patterns.}
Doji: $|O - C| < 0.1\%$ of the range $H - L$. Hammer: small body
($|O - C| < 0.3 (H - L)$) and long lower shadow
($\min(O, C) - L > 2 |O - C|$). Both are one-line tests inside the
\func{Ohlcv} struct's accessor methods.

\paragraph{2. Volume Analysis.}
\func{volume\_ma} is a simple \func{sma} applied to the volume field.
A day with volume $> 2 \times \text{MA}_N$ is flagged as high-volume.
On the bundled \texttt{stock\_prices.csv}, roughly 5--8\% of days
exceed 2x; these often coincide with earnings or macro releases.

\paragraph{3. Bollinger Bands.}
Upper $= \text{SMA}_N + 2 \sigma_N$, lower $= \text{SMA}_N - 2
\sigma_N$, where $\sigma_N$ is the rolling standard deviation (use
\func{rolling\_std\_dev} from \crate{quant-core}). Prices touch the
upper band on $\sim 5\%$ of days under a normal approximation,
confirming the bands are a 2-sigma envelope.

\paragraph{4. MACD.}
$\text{MACD} = \text{EMA}_{12} - \text{EMA}_{26}$,
$\text{signal} = \text{EMA}_9(\text{MACD})$. A buy signal fires when
MACD crosses above its signal line. The indicator is a
momentum oscillator; it lags the price turn because EMAs are
backward-looking.

\paragraph{5. Backtest.}
Buy at the close on a golden cross, sell at the close on a death
cross. On the synthetic series, the strategy is profitable only if
the trend persists longer than the lag; in choppy markets it
whipsaws and loses. The lesson: SMA crossover is a trend filter,
not a predictor---it works in regimes and fails in transitions.

\section*{Chapter 4: Returns and Volatility}

\paragraph{1. Return Conversion.}
$r^{\text{log}} = \ln(1 + r^{\text{simple}})$ and
$r^{\text{simple}} = e^{r^{\text{log}}} - 1$. A property test
generates random simple returns, converts to log and back, and
asserts the round-trip error is below $10^{-15}$. The identity is
exact in real arithmetic; floating-point round-off is the only
source of error.

\paragraph{2. Annualisation Sanity Check.}
Resample daily returns to weekly and monthly. Annualised volatility
is $\sigma_{\text{daily}} \sqrt{252} \approx
\sigma_{\text{weekly}} \sqrt{52} \approx
\sigma_{\text{monthly}} \sqrt{12}$. The three estimates agree to
within sampling noise; the discrepancy is $O(1/\sqrt{T})$ where $T$
is the number of observations at each frequency.

\paragraph{3. Calmar Ratio.}
$\text{Calmar} = \text{annualised return} / |\text{max drawdown}|$.
On a synthetic strategy with 12\% annual return and 15\% max
drawdown, Calmar $= 0.8$. Sharpe penalises all volatility; Calmar
penalises only the worst drawdown. A strategy with high Sharpe but
deep drawdown has a low Calmar---common in trend-following.

\paragraph{4. Rolling Sharpe.}
Windowed Sharpe $= \frac{\bar{r}_N - r_f}{\sigma_N} \sqrt{m}$. It
dips during volatile periods (e.g., a 2-sigma shock raises $\sigma_N$
immediately but the mean responds slowly). The rolling chart
visualises regime changes that the full-sample Sharpe hides.

\paragraph{5. Drawdown with Recovery.}
Average drawdown $= \frac{1}{T} \sum_t \min(0, \text{DD}_t)$.
Proportion in drawdown $= \frac{1}{T} \sum_t
\mathbb{1}[\text{DD}_t < 0]$. A buy-and-hold equity curve spends
60--80\% of its time in drawdown---the ``tyranny of drawdowns'' that
motivates risk-managed strategies.

\section*{Chapter 5: Backtesting}

\paragraph{1. Cost Sensitivity.}
For $\tau \in \{0, 0.001, 0.0025, 0.005, 0.01\}$, total return
declines roughly linearly in $\tau$ because turnover is roughly
constant. The strategy stops being profitable when
$\tau \times \text{turnover} > \text{gross edge}$. For the 10/30
crossover on the synthetic series, this happens near
$\tau \approx 0.005$ (50 bps round-trip).

\paragraph{2. Parameter Sensitivity.}
The $(S, L)$ that maximises Sharpe is not the one that maximises
total return: a faster crossover (small $S, L$) has higher gross
return but higher turnover, so net return (after costs) is lower.
The Sharpe-maximising configuration is typically $(10, 50)$ or
$(20, 100)$---slow enough to filter noise, fast enough to catch
trends.

\paragraph{3. EMA Crossover.}
EMA crossover has lower turnover than SMA because EMAs decay
smoothly rather than dropping observations. On the same data, the
EMA strategy typically trades 20--30\% less and has a comparable or
slightly higher Sharpe. The trade-off: EMAs react faster but are
more sensitive to outliers.

\paragraph{4. Position Sizing.}
With fraction $f = 0.5$, the equity curve is flatter (lower return,
lower volatility, shallower drawdowns). The trade-off is
return-for-risk: $f = 0.5$ roughly halves both the annual return
and the max drawdown. Full-Kelly ($f = f^*$) maximises log growth;
fractional Kelly trades growth for safety.

\paragraph{5. Walk-Forward.}
Tune $(S, L)$ on the first 180 bars; freeze; evaluate on the last
72. The out-of-sample Sharpe is typically 30--60\% of the in-sample
Sharpe---a WFE of 0.3--0.6. A WFE below 0.3 signals overfitting; the
in-sample optimum exploited noise.

\section*{Chapter 6: Statistical Foundations}

\paragraph{1. Sample vs Population Variance.}
For $\{1,2,3,4,5\}$: population variance $= 2.0$ (denominator $n=5$),
sample variance $= 2.5$ (denominator $n-1=4$). The bias is
$(n-1)/n = 0.8$, a 20\% underestimate. It matters for small samples
and for any statistic that scales with variance (Sharpe, t-stat).
Always use $n-1$ for sample-based estimators.

\paragraph{2. Skewness Sign.}
Negative skew: mix a normal with a few large negative draws
(e.g., $x = N(0,1)$ with probability 0.9, $x = -5$ with probability
0.1). Positive skew: the mirror. \func{skewness} returns the correct
sign. Equity returns typically have \emph{negative} skew---large
drops are more frequent than large jumps.

\paragraph{3. RNG Period.}
A test asserts the first 1000 outputs from seed 42 are distinct.
``Distinct'' is necessary (a PRNG that repeats early is broken) but
not sufficient (a PRNG can be distinct yet correlated---e.g., a
counter). The full statistical battery (\texttt{TestU01},
\texttt{DieHarder}) checks uniformity, independence, and
higher-order structure.

\paragraph{4. GBM Convergence.}
With $S_0 = 100$, $\mu = 0.05$, $\sigma = 0.2$, $T = 1$, the
analytical $\E[S_T] = 100 \cdot e^{0.05} = 105.13$. A 10\,000-path
Monte Carlo gives $\hat{\E}[S_T] \approx 105.1 \pm 0.2$ (the SE
scales as $\sigma \sqrt{T} / \sqrt{N}$). The estimate is within
one SE of the truth---the small discrepancy is sampling noise, not
bias.

\paragraph{5. Jump-Diffusion.}
Add $J \sim N(0, 0.05^2)$ with intensity $\lambda$. As $\lambda$
grows from 0 to 1 (one jump per year on average), the excess
kurtosis of terminal prices rises from 0 to $\sim 3$--$5$. Jumps
inject fat tails that GBM alone cannot reproduce---the basis of
the Merton model.

\section*{Chapter 7: Time Series}

\paragraph{1. Partial Autorrelation.}
The PACF via Yule-Walker solves the linear system
$\mathbf{R} \boldsymbol{\phi} = \mathbf{r}$ where $R_{ij} =
\rho_{|i-j|}$ and $r_i = \rho_i$. For AR(2) with $\phi_1 = 0.5$,
$\phi_2 = -0.3$, PACF is $\sim 0.5$ at lag 1, $\sim -0.3$ at lag 2,
and $\sim 0$ beyond---the cutoff at lag $p = 2$ identifies the
order.

\paragraph{2. KPSS Test.}
KPSS tests $H_0$: stationary (opposite of ADF). On white noise, both
ADF and KPSS agree (stationary). On a random walk, ADF fails to
reject (non-stationary) and KPSS rejects (non-stationary). On
fractionally differenced series with $d < d^\star$, the two can
disagree: ADF may reject (stationary by unit-root criterion) while
KPSS rejects (non-stationary by trend criterion). The disagreement
signals a long-memory process---neither pure stationary nor pure
unit root.

\paragraph{3. Real Returns.}
For $\sigma \in \{0.1, 0.2, 0.4\}$, the minimum $d^\star$ for
stationarity is approximately $0.35$--$0.45$ and is largely
insensitive to $\sigma$. The reason: $d^\star$ depends on the
memory structure (how slowly the ACF decays), not the scale. Higher
$\sigma$ adds noise but does not change the differencing order
needed to remove the unit root.

\paragraph{4. Expanding-Window Frac-Diff.}
The expanding window uses all available history for early points,
truncating the weight sequence at the available length rather than
dropping the first $K-1$ points. The trade-off: more data for early
points (no warm-up loss) but the early values are biased because
they use fewer terms of the infinite series. Fixed-width is
consistent (every point uses the same number of terms); expanding
is efficient (no data loss) but heterogeneous.

\paragraph{5. Lag Selection for ADF.}
AIC $= n \ln(\hat{\sigma}^2) + 2k$ where $k$ is the number of lags.
Sweep $k$ from 0 to 10, pick the AIC-minimising $k$, then run ADF.
The stationarity conclusion rarely changes for well-behaved series;
it can change for borderline series where the lag structure affects
the residual autocorrelation. AIC-based selection is more
principled than a fixed lag.

\section*{Chapter 8: Volatility Models}

\paragraph{1. EGARCH.}
$\log \sigma_t^2 = \omega + \alpha g(z_{t-1}) + \beta \log
\sigma_{t-1}^2$ where $g(z) = \theta z + \gamma(|z| - \E|z|)$. The
$\theta$ term captures the leverage effect: negative returns
($z < 0$) raise $\log \sigma^2$ more than positive ones. On equity
returns, EGARCH's log-likelihood typically exceeds GARCH(1,1) by
5--15 points, with $\theta < 0$ significantly.

\paragraph{2. GJR-GARCH.}
$\sigma_t^2 = \omega + \alpha r_{t-1}^2 + \gamma
\mathbb{1}_{r_{t-1}<0} r_{t-1}^2 + \beta \sigma_{t-1}^2$. On equity
index returns, $\gamma$ is typically $0.03$--$0.10$ and
significant, confirming the leverage effect. The effective
coefficient on negative shocks is $\alpha + \gamma$, larger than
$\alpha$ alone.

\paragraph{3. Forecast Evaluation.}
Fit on 80\%, forecast the conditional variance on the remaining 20\%.
The GARCH forecast converges to the long-run variance
$\sigma_\infty^2 = \omega / (1 - \alpha - \beta)$ as the horizon
grows. The mean squared error of the variance forecast is typically
20--40\% lower than the EWMA forecast, because GARCH adapts to the
estimated persistence while EWMA uses a fixed $\lambda = 0.94$.

\paragraph{4. Student-$t$ Innovations.}
Replace the Gaussian density with the Student-$t$ density with $\nu$
degrees of freedom. On heavy-tailed returns, the fitted $\nu$ is
typically 4--7, and the $t$-GARCH log-likelihood exceeds the Gaussian
GARCH by 10--30 points. The improvement is largest in the tails,
which the Gaussian model underweights.

\paragraph{5. Long-Memory FIGARCH.}
FIGARCH($1,d,1$) replaces the geometric decay with a hyperbolic
$(1-L)^d$ filter on the variance recursion. The half-life of
shocks grows from $\sim 10$ days (GARCH(1,1)) to $\sim 30$--$50$
days (FIGARCH), reflecting the slow decay of volatility
autocorrelation observed in high-frequency data.

\section*{Chapter 9: Stochastic Processes}

\paragraph{1. Control Variates.}
$\hat{C}_{\text{cv}} = \hat{C} - \beta(\hat{M} - \E[S_T])$ with
$\beta = \Cov(\text{payoff}, S_T) / \Var(S_T)$. For a call option,
$\beta \approx e^{-rT} N(d_1)$ (the delta). The standard error drops
by a factor of $\sqrt{1 - \rho^2}$ where $\rho$ is the
correlation between payoff and $S_T$---typically 70--90\% for
vanilla calls, giving a 2--3x SE reduction.

\paragraph{2. Asian Option.}
Time-step the GBM over $n = 252$ daily closes, compute $\bar S =
\frac{1}{n} \sum S_{t_k}$, payoff $\max(\bar S - K, 0)$. The Asian
price is below the European price (averaging reduces variance, so
the option is cheaper). The SE is comparable to the European for
the same $N$, because the variance reduction from averaging is in
the payoff, not the estimator.

\paragraph{3. Euler vs Exact.}
Euler discretisation $S_{t+\Delta t} = S_t(1 + \mu \Delta t + \sigma
\sqrt{\Delta t} Z)$ has bias $O(\Delta t)$ relative to the exact
$S_T = S_0 \exp((\mu - \tfrac{1}{2}\sigma^2)T + \sigma W_T)$. At
$\Delta t = 1/252$, the bias in the mean is $\sim 0.01\%$---negligible
for pricing but visible in high-precision convergence studies.

\paragraph{4. Confidence Intervals.}
For $N = 10\,000$, the 95\% CI is $\hat{C} \pm 1.96 \cdot
\text{SE}$. Repeating 1000 times, the fraction of CIs covering the
BS price is $\sim 94$--$96\%$, consistent with the nominal 95\%.
Coverage below 93\% signals a bug (e.g., wrong discounting); coverage
above 97\% suggests the SE is overestimated.

\section*{Chapter 10: Options Pricing}

\paragraph{1. Black-Scholes with Dividends.}
Replace $S_0$ with $S_0 e^{-qT}$. The call becomes $C = S_0 e^{-qT}
N(d_1) - K e^{-rT} N(d_2)$ with $d_1 = \frac{\ln(S_0/K) + (r - q +
\tfrac{1}{2}\sigma^2)T}{\sigma\sqrt{T}}$. Delta $= e^{-qT} N(d_1)$
(the dividend yield discounts the delta). The forward price
$F = S_0 e^{(r-q)T}$ replaces $S_0 e^{rT}$ in put-call parity.

\paragraph{2. Local Volatility.}
$\sigma(S, t) = \sigma_0 + \alpha (S - S_0)^2$ creates a smile: ATM
has vol $\sigma_0$, OTM and ITM have higher vol. Explicit Euler on
a $(S, t)$ grid with $\Delta S = 1$, $\Delta t = 1/252$ prices the
option. The IV surface recovered by inverting BS on the grid prices
is convex in moneyness, matching the local-vol shape.

\paragraph{3. Heston Model.}
$dv_t = \kappa(\theta - v_t) dt + \xi \sqrt{v_t} dW^{(v)}_t$ with
$\Cov(dW, dW^{(v)}) = \rho \, dt$. Monte Carlo with full truncation
($v_t = \max(v_{t-1}, 0)$) generates call prices across a strip of
strikes. The IV smile is skewed: negative $\rho$ produces a
downward-sloping skew (puts have higher IV), matching equity
markets.

\paragraph{4. Vega-Clipped Bisection.}
For $S_0 = 1000$, $K = 10$, the call price equals the intrinsic
$990 \cdot e^{-rT}$ to machine precision. Vega is essentially zero
(dividing by $d_2 \to \infty$), so Newton's method diverges. The fix:
detect $\text{Vega} < \epsilon$ and return
\texttt{Ok(sigma\_min)} with a warning flag, because the IV is
genuinely unrecoverable---any vol gives the same price.

\paragraph{5. Smile Fitting.}
Fit $\sigma(K) = \beta_0 + \beta_1 \ln(K/S_0) + \beta_2
(\ln(K/S_0))^2$ by OLS. On SPX option data, the residual is
typically 0.1--0.5 vol points, $\beta_1$ is negative (the skew),
and $\beta_2$ is positive (the convexity). The skew coefficient
$\beta_1$ quantifies how much the market prices downside risk.

\paragraph{6. Newton Convergence.}
Instrument \func{implied\_vol} to print $|C(\sigma_n) -
C_{\text{mkt}}|$ at each step. For ATM options, the error roughly
squares each iteration (quadratic convergence) until it hits
machine precision in 4--6 steps. For deep ITM/OTM, Newton slows to
linear (the derivative is tiny), and the bisection fallback takes
over, converging in $\sim 10$--$20$ iterations.

\section*{Chapter 11: Portfolio Optimization}

\paragraph{1. Correlated Assets.}
For $\rho = 0.5$, the frontier bends left (diversification helps
but less than independence). For $\rho = -0.5$, it bends further
left. For $\rho = -1$, the frontier touches the $\sigma = 0$ axis:
the zero-variance portfolio has weight $w_A = \sigma_B / (\sigma_A +
\sigma_B)$, giving $\sigma_p = 0$ exactly. This is the only case
where risk can be fully eliminated.

\paragraph{2. Three-Asset Frontier.}
With $\mu_C = 0.08$, $\sigma_C = 0.15$, the global minimum-variance
portfolio mixes all three assets to minimise $\mathbf{w}^\top
\boldsymbol{\Sigma} \mathbf{w}$. The tangency Sharpe is higher than
the two-asset case because the third asset adds a diversification
opportunity. The gain is largest when the new asset has low
correlation with the existing two.

\paragraph{3. Long-Only Constraint.}
Projected gradient descent: take a gradient step, then project onto
the probability simplex ($w_i \ge 0$, $\sum w_i = 1$) by sorting and
thresholding. For a high target $\mu_\star = 0.15$, the unconstrained
solution shorts the low-return asset; the constrained solution
allocates 100\% to the high-return asset, with higher variance. The
constraint widens the frontier.

\paragraph{4. Black-Litterman.}
Prior $\mu \propto \Sigma \mathbf{w}_{\text{mkt}}$ (reverse-optimised
from market weights). Posterior $\mu_{\text{BL}} = [(\tau\Sigma)^{-1}
+ \Omega^{-1}]^{-1} [(\tau\Sigma)^{-1}\mu + \Omega^{-1}\mathbf{q}]$
where $\mathbf{q}$ is the view and $\Omega$ the view confidence. The
posterior tangency tilts toward the view; the prior tangency is the
market portfolio. The shrinkage parameter $\tau$ controls how much
the view moves the estimate.

\paragraph{5. Rolling CAPM.}
Rolling 60-day beta $= \Cov(r_i, r_m) / \Var(r_m)$ computed over a
sliding window. Betas are time-varying: a defensive stock's beta
drops in bear markets (investors flock to safety), a cyclical
stock's beta rises. The plot visualises the regime change---a
constant beta assumption is a first-order approximation.

\paragraph{6. CVaR Optimisation.}
CVaR$_{95\%} = \E[-r_p \mid -r_p > \text{VaR}_{95\%}]$. Minimising
CVaR over the two-asset simplex by grid search gives a portfolio
that is more concentrated in the lower-variance asset than the
Markowitz minimum-variance portfolio. The two differ when the return
distribution is asymmetric---CVaR penalises tail shape, variance
penalises spread.

\section*{Chapter 12: Factor Models}

\paragraph{1. Four- and Five-Factor Models.}
Add momentum (UMD) and quality (e.g., gross profitability). The
$R^2$ typically rises by 2--8 percentage points over the 3-factor
model, with momentum contributing more for individual stocks and
quality more for portfolios. The incremental gain diminishes as more
factors are added---diminishing returns to factor zoo.

\paragraph{2. Jacobi Eigenvalue.}
The Jacobi algorithm rotates pairs of rows/columns to zero out
off-diagonal entries. For a $5 \times 5$ symmetric matrix, it
converges in 5--15 sweeps, each $O(n^2)$ rotations. Compared to the
power method, Jacobi computes all eigenvalues simultaneously (no
deflation error) but is slower for large $n$. For $n = 5$ the two
agree to machine precision.

\paragraph{3. PCA on Real Equity Returns.}
Run PCA on 10 stocks, 252 days. The first PC is typically a ``market
factor'' (all-positive loadings, $\sim 30$--$50\%$ of variance). The
second and third PCs capture sector tilts (e.g., growth vs value,
financials vs tech). The economic interpretation is not
automatic---it requires labelling the loadings by sector.

\paragraph{4. Rolling 60-Day Beta.}
$\beta_t = \Cov_{60}(r_i, r_m) / \Var_{60}(r_m)$, plotted over time.
Betas of defensive stocks (utilities) are $< 1$ and stable; betas
of cyclical stocks (tech, finance) are $> 1$ and volatile. The
rolling chart shows when a stock's sensitivity to the market
changes---useful for regime-aware risk management.

\paragraph{5. Market-Neutral Long-Short.}
Long the high-SMB quintile, short the low-SMB quintile, scaled so
$\beta_{\text{Mkt}} = 0$. The portfolio's variance decomposes into
systematic (SMB-driven) and idiosyncratic (stock-specific) parts.
Diversification (many stocks per leg) eliminates the idiosyncratic
part, leaving the systematic SMB exposure---a ``pure size factor''
portfolio.

\section*{Chapter 13: Market Microstructure}

\paragraph{1. Order Book Builder.}
\begin{lstlisting}[language=Rust]
fn build_book(orders: Vec<(Side, i64, u64, u64)>) -> OrderBook {
    let mut book = OrderBook::new(1);
    for (side, price, qty, ts) in orders {
        book.add_order(Order { id: ts, side, price, quantity: qty, timestamp: ts }).unwrap();
    }
    book
}
\end{lstlisting}
Assert \func{best\_bid} and \func{best\_ask} return the expected
prices and \func{total\_volume} sums the quantities.

\paragraph{2. Liquidity Sweep.}
Build 5 ask levels totalling 1000 shares. A market buy of 750
sweeps the first $\sim 3$--$4$ levels. The VWAP of the fills is
$\le$ the volume-weighted ask price of all swept levels (equality
when the order exactly exhausts a level). The test confirms
price-time priority: earlier orders fill first.

\paragraph{3. OFI and Returns.}
Generate snapshots where bid quantity rises with ``buy pressure''
and ask quantity falls. OFI $= \Delta q^{\text{bid}} - \Delta
q^{\text{ask}}$ is positive when buy pressure rises. Mid-price
returns are also positive. The correlation $\rho(\text{OFI},
r_{\text{mid}})$ is typically 0.3--0.6, confirming OFI is a
signed predictor of short-term returns.

\paragraph{4. Impact Scaling.}
\func{sqrt\_impact}$(\sigma, Q, V) = \sigma \sqrt{Q/V}$. Plotting
impact vs $Q$ on a log-log scale gives a line of slope $1/2$,
confirming the square-root law. Doubling $Q$ multiplies impact by
$\sqrt{2} \approx 1.41$, not by 2---the concavity is the empirical
regularity that motivates the Almgren-Chriss model.

\paragraph{5. Net Frontier.}
Per-rebalance cost $= \func{execution\_cost}(Q = 0.5 V_{\text{daily}},
V_{\text{daily}}, \sigma, \text{spread})$. At $50\%$ turnover, the
cost is $\sim 5$--$15$ bps. Subtract from the expected return: the
net efficient frontier shifts down by the cost times turnover. The
tangency Sharpe drops by $\sim 0.05$--$0.15$, depending on the
spread and impact. The lesson: gross frontier optimisation
overstates net performance.

\section*{Chapter 14: AFML Backtesting}

\paragraph{1. Custom Barriers.}
Set the upper barrier to $\phi_t = 2 \sigma_t$ and the lower to
$\psi_t = 2 \sigma_t$ where $\sigma_t$ is the rolling 20-day
volatility. On a noisy series, the barriers widen, fewer events
hit the profit-taking label (more time-outs); on a calm series,
the barriers narrow, more events hit the barriers. The label
distribution shifts from $\sim 20\%$ up to $\sim 10\%$ up as
volatility rises---volatility-scaled barriers adapt to the regime.

\paragraph{2. Embargo Calibration.}
Compute the ACF of FFD returns. The embargo $\rho$ is the first lag
where $|\rho_k| < 1.96 / \sqrt{T}$ (the 5\% significance band). For
daily FFD returns, $\rho$ is typically 1--5 days. Without embargo,
the out-of-sample Sharpe is inflated by $\sim 5$--$15\%$ due to
leakage; with embargo, the Sharpe drops to its honest value.

\paragraph{3. Volatility-Scaled Kelly.}
$f_t = f^* \cdot \bar{\sigma} / \sigma_t$ where $\bar{\sigma}$ is
the long-run volatility. In high-vol periods, $f_t$ shrinks; in
low-vol periods, it grows. The drawdown distribution tightens
because the largest losses (which occur in high-vol periods) are
scaled down. The trade-off: the strategy grows more slowly in calm
periods, but survives turbulent ones.

\paragraph{4. Meta-Labeling.}
Primary: sign of SMA crossover. Label with triple barrier. Train a
logistic regression on the labelled events to predict correctness
(meta-model). The meta-filtered strategy takes the primary signal
only when the meta-model predicts ``correct'' (probability $> 0.5$).
The meta-filtered Sharpe is typically 20--50\% higher than the
primary alone, with fewer trades---the meta-model filters out the
primary's low-confidence signals.